\section{Условие}

Вариант 2. Родительский процесс считывает числа с плавающей запятой.
Дочерний процесс вычисляет их сумму, записывает результат в указанный пользователем
файл и возвращает сумму родителю. Обмен выполняется через отображаемый файл в Linux.

\section{Метод решения}

Оба процесса отображают один временный файл с флагом \texttt{MAP\_SHARED}.
В начале отображения находится структура \texttt{shared\_data}. Поле \texttt{count}
хранит количество чисел, \texttt{sum} содержит итог, а массив \texttt{values[]}
содержит введённые числа.

\begin{enumerate}
  \item \texttt{parent} создаёт временный файл через \texttt{mkstemp()} и удаляет
    его имя вызовом \texttt{unlink()}. Файл остаётся доступен по открытому дескриптору.
    Родитель задаёт начальный размер через \texttt{ftruncate()} и отображает файл
    вызовом \texttt{mmap()}. При заполнении массива он увеличивает файл через
    \texttt{ftruncate()}, а отображение расширяет через \texttt{mremap()}.
  \item После пустой строки или конца ввода родитель запускает \texttt{child} через
    \texttt{fork()} и \texttt{execl()}. Дескриптор файла сохраняется при запуске
    новой программы. Потомок узнаёт размер файла через \texttt{fstat()} и создаёт
    собственное отображение через \texttt{mmap()}.
  \item Потомок суммирует элементы \texttt{values[]}, сохраняет результат в
    выходной файл и записывает сумму в поле \texttt{sum}. Родитель читает это поле
    после успешного завершения потомка, которое проверяет через \texttt{waitpid()}.
\end{enumerate}

Родитель заканчивает запись чисел до запуска потомка, а результат читает после
его завершения. Такой порядок исключает одновременный доступ к полям, которые
один процесс записывает, а другой читает.

При сбоях программа печатает диагностику в \texttt{stderr} и возвращает код 1.
Оба процесса снимают отображение через \texttt{munmap()} и закрывают дескриптор
через \texttt{close()}.

\section{Вывод}

В ходе работы освоены алгоритм обмена данными между процессами через отображаемый
файл, динамическое расширение общего буфера и последовательное суммирование
массива в дочернем процессе.

В лабораторной работе №~2 данные передавались через каналы как поток байтов
с помощью операций чтения и записи. В этой реализации оба процесса обращаются
к полям одной структуры через отображения общего файла. Родитель записывает
числа до запуска потомка и читает сумму после его завершения, поэтому отдельные
семафоры для обмена не понадобились.
